\subsection{Glyph: \glyph{Phenotype}}
\label{sec:phenotype}

A biochemical network can generate phenotypes or affect biological processes.
Such processes can take place at different levels and are independent of the biochemical network itself.
To represent these processes in a map, \PD defines the \glyph{phenotype} glyph.

\begin{glyphDescription}

\glyphSboTerm
SBO:0000358 ! phenotype


\glyphIncoming
One or more \glyph{modulation} arcs (\sect{modulations}).



\glyphOutgoing
None.


\glyphContainer
A \glyph{phenotype} is represented by an elongated hexagonal shape, as shown in \fig{phenotype}.

\glyphLabel
A \glyph{phenotype} is identified by a label that is  a string of characters that may be distributed on several lines to improve readability.
The centre of the label must be placed on the centre of the shape.
The label may extend outside of the shape.

\glyphAux 
None.

\end{glyphDescription}

\begin{figure}[H]
  \centering
  \includegraphics{images/build/phenotype.pdf}
  \caption{The \PD glyph for \glyph{phenotype}.}
  \label{fig:phenotype}
\end{figure}

% The following is for [X]Emacs users.   Please leave in place.
% Local Variables:
% TeX-master: "../sbgn_PD-level1"
% End:
